%% 
%% ACS project dissertation template. 
%% 
%% Currently designed for printing two-sided, but if you prefer to 
%% print single-sided just remove ",twoside,openright" from the 
%% \documentclass[] line below. 
%%
%%
%%   SMH, May 2010. 

\documentclass[a4paper,12pt,twoside,openright]{report}
\usepackage{graphicx}

%%
%% EDIT THE BELOW TO CUSTOMIZE
%%

\def\authorname{Martin Marinov\xspace}
\def\authorcollege{St Edmund's College\xspace}
\def\authoremail{mtm46@cam.ac.uk}
\def\dissertationtitle{Remote video eavesdropping using a software-defined radio platform}
\def\wordcount{??,???}


\usepackage{epsfig,graphicx,parskip,setspace,tabularx,xspace,epstopdf,amssymb,caption,subcaption} 
\graphicspath{ {./images/} }

%% START OF DOCUMENT
\begin{document}


%% FRONTMATTER (TITLE PAGE, DECLARATION, ABSTRACT, ETC) 
\pagestyle{empty}
\singlespacing
% title page information
\begin{titlepage} 

\begin{center}
\noindent
\huge
\dissertationtitle \\
\vspace*{\stretch{1}}
\end{center}

\begin{center}
\noindent
\huge
\authorname \\
\Large
\authorcollege      \\[24pt]
\includegraphics{CUni3.eps}
\end{center}

\vspace{24pt} 

\begin{center}
\noindent
\large
{\it A dissertation submitted to the University of Cambridge \\ 
in partial fulfilment of the requirements for the degree of \\ 
Master of Philosophy in Advanced Computer Science} 
\vspace*{\stretch{1}}
\end{center}

\begin{center}
\noindent
University of Cambridge \\
Computer Laboratory     \\
William Gates Building  \\
15 JJ Thomson Avenue    \\
Cambridge CB3 0FD       \\
{\sc United Kingdom}    \\
\end{center}

\begin{center}
\noindent
Email: \authoremail \\
\end{center}

\begin{center}
\noindent
\today
\end{center}

\end{titlepage} 

\newpage
\vspace*{\fill}

\onehalfspacing
\newpage
{\Huge \bf Declaration}

\vspace{24pt} 

I \authorname of \authorcollege, being a candidate for the M.Phil in
Advanced Computer Science, hereby declare that this report and the
work described in it are my own work, unaided except as may be
specified below, and that the report does not contain material that
has already been used to any substantial extent for a comparable
purpose.

\vspace{24pt}
Total word count: \wordcount

\vspace{60pt}
\textbf{Signed}: 

\vspace{12pt}
\textbf{Date}:


\vfill

This dissertation is copyright \copyright 2014 \authorname. 
\\
All trademarks used in this dissertation are hereby acknowledged.



\newpage
\vspace*{\fill}

\singlespacing
\newpage
{\Huge \bf Abstract}
\vspace{24pt} 


This dissertation presents a software toolkit for remotely eavesdropping video monitors using a Software Defined Radio (SDR) receiver. It exploits compromising emanations from cables carrying video signals.

Raster video is usually transmitted one line of pixels at a time, encoded as a varying current. This generates an electromagnetic wave that can be picked up by an SDR receiver. The software maps the received field strength of a pixel to a gray-scale shade in real-time. This forms a false colour estimate of the original video signal.

The toolkit uses unmodified off-the-shelf hardware which lowers the costs and increases mobility compared to existing solutions. It allows for additional post-processing which improves the signal-to-noise ratio. The attacker does not need to have prior knowledge about the target video display. All parameters such as resolution and refresh rate are estimated with the aid of the software. 

The software consists of a library written in C, a collection of plug-ins for various Software Define Radio (SDR) front-ends and a Java based Graphical User Interface (GUI). It is a multi-platform application, with all native libraries pre-compiled and packed into a single Java jar file.


\newpage
\vspace*{\fill}


\pagenumbering{roman}
\setcounter{page}{0}
\pagestyle{plain}
\tableofcontents
%% \listoffigures
%% \listoftables

\onehalfspacing

%% START OF MAIN TEXT 

\chapter{Introduction}
\pagenumbering{arabic} 
\setcounter{page}{1} 

\section{Related Work} 

[Related work and critique]

\section{Availability} 
[Github repo url]

\cite{kuhn2003compromising}

\chapter{Background}

\section{Electronic Emanations} 

\section{IQ sampling} 

\chapter{Methodology} 

\section{Analogue Video Signals}

Almost all contemporary video monitors are raster based. The image is transferred from the video controller to the display in scanlines that occur at a specific rate. Each of these scanlines contains a number of pixels which are continuously encoded as a time varying signal. This signal is generated internally in the video controller by an oscillator which runs at the pixel clock rate. The signal is thereafter multiplied by the pixel intensities at the specific time.

Let's assume that the signal started transmitting at time $t=0$, and each video frame contains $h$ scanlines, each of which contains $w$ pixels. The frequency with which frames are being generated is $f_{f}$ frames per second. Therefore the duration of transmission of each individual pixel is \begin{equation}
\label{eq:tdelta_definition}
T_{\Delta}=\frac{1}{w h f_{f}}
\end{equation}
Furthermore at time $t$ frame number $n_{t}$ started transmission where 
\begin{equation}
n_{t}=\lfloor t f_{f} \rfloor
\end{equation}
Therefore 
if we assume that the top left corner of a frame has coordinates $(0, 0)$, a pixel at position $(x, y)$ in frame $n_{t}$ will start to be transmitted at time
\begin{equation}
T_{(x,y)}= (y w + x + n_{t} w h) T_{\Delta}
\end{equation}
and will finish transmission before time 
\begin{equation}
T_{(x+1,y)} = T_{(x,y)} + T_{\Delta} =(y w + x + 1 + n_{t} w h) T_{\Delta}
\end{equation}
at which the next pixel will start transmitting.

In practice $w$ and $h$ are determined by the screen resolution and $f_{f}$ is simply the screen refresh rate. For a typical screen resolution of $width \times height$, it is true that $w \geq width$ and $h \geq height$. The reason is that video signals tend to have additional blanking intervals. This means more pixels are transmitted than what is in the active video region. This gives opportunity for the receiving monitor to synchronise its internal clock, calibrate its colour levels or in case of CRT, allow enough time for the electron beam to return to the beginning of the next line on the screen. The synchronisation timings for personal computers have been standardised by Video Electronics Standards Association (VESA) \cite{vesa}.

In order to decode an individual pixel, the receiving monitor has its own internal oscillator. It locks it to the pixel rate of the incoming signal either via an external clock source or using the blanking intervals. Once it receives the signal for an individual pixel, its amplitude (or binary content in case of digital signal) will correspond to the intensity. This allows the monitor to display the video in real time. If multiple colours are desired, they can be transmitted separately on different wires in the same fashion.

\section{Generated Radio Wave}

\subsection{Signal Equation}
Let's assume the discrete pixels in a video signal have intensities $v_{i} (i \in \mathbb{Z})$ and are being transmitted for a duration of $T_{\Delta}=\frac{1}{w h f_{f}}$ (from \ref{eq:tdelta_definition}). Let's also assume that the shape of the pixel is $p(t)$ where $p(t)=0$ for $|t| > \frac{T_{\Delta}}{2}$. We know that pixel $i$ starts transmitting at time $t_{i}=i T_{\Delta}$ (if we assume that that pixel 0 was transmitted at $t=0$). Also the amplitude of the signal of the transmitted pixel is linearly dependant on its intensity $v_{i}$. Then the resulting signal between time $t_{i}$ and $t_{i+1}$ will have the form $v_{i} p(t-i T_{\Delta})$. To generalise for all $i$, the resulting actual transmitted radio wave in the time domain will have the form 
\begin{equation}
\label{eq:video_signal}
\tilde{v}(t) = \sum\limits_{i=-\infty}^{\infty} v_{i} p(t-i T_{\Delta})
\end{equation}
Which is a continuous function. Equation \ref{eq:video_signal} can be rewritten as a convolution with a Dirac delta function
\begin{equation}
\label{eq:vasconv}
\tilde{v}(t) = p(t) \ast \left( \sum\limits_{i=-\infty}^{\infty} v_{i} \cdot \delta(t-i T_{\Delta}) \right) = p(t) \ast \hat{v}(t)
\end{equation}
Where we have put
\begin{equation}
\label{eq:vhattdef}
\hat{v}(t) = \sum\limits_{i=-\infty}^{\infty} v_{i} \cdot \delta(t-i T_{\Delta})
\end{equation}
Note that this results in infinitesimally short (in time) spikes at values exact integer multiples of $T_{\Delta}$ with amplitudes weighted to the pixel intensities at that time. We can view \ref{eq:vasconv} as the pixel shape being repeated at intervals of $T_{\Delta}$ modulated by $v_{i}$. 

\subsection{Sampling}
We can notice that \ref{eq:vhattdef} looks like the result of a continuous signal being sampled at discrete points in time $T_{\Delta}$ apart. If we call this continuous signal $v(t)$, then we can rewrite the equation so that it reads
\begin{equation}
\hat{v}(t) = \sum\limits_{i=-\infty}^{\infty} v(t) \cdot \delta(t-i T_{\Delta}) = v(t) \sum\limits_{i=-\infty}^{\infty} \delta(t-i T_{\Delta})
\end{equation}
The requirement is that if $v(t)$ is sampled at discrete intervals $i T_{\Delta}$, it should be equal to the corresponding pixel values i.e. $v(0)=v_{0}$, $v(T_{\Delta})=v_{1}$, $v(2 T_{\Delta})=v_{2}$, etc. This would mean that we can precisely obtain the samples from the continuous function by applying the series of Dirac delta functions. However, now we would also like to be able to obtain the continuous signal by only having the sampled values $v_{i}$. This means that we need a unique perfect reconstruction i.e. to be able to obtain all $v_{i}$ from $v(x)$ and $v(x)$ from all $v_{i}$. This condition is satisfied by the Whittaker--Shannon interpolation formula which can give us the condition for perfect reconstruction for digitally sampled signals. It states that the continuous signal could be uniquely reconstructed using $sinc$ interpolation. If we apply it to our $v_{i}$ samples, it yields
\begin{equation}
v(t) = \sum\limits_{k=-\infty}^{\infty} v_{k} \cdot sinc ( \pi \left( \frac{t}{T_{\Delta}} - k \right) )
\end{equation}
We can verify that when $t=i T_{\Delta}$, then $sinc(\pi (i-k) )$ will always yield $0$ except for when $i = k$ in which case it will be $1$. Therefore we have proved that $v(i T_{\Delta}) = v_{i}$.

\subsection{Repetitions}
Now let's revise some mathematical identities. According to the convolution theorem
\begin{equation}
\mathcal{F} \left\{ g \cdot h \right\} = \mathcal{F} \left\{ g \right\} \ast  \mathcal{F} \left\{ h \right\}
\end{equation}
Which reads: the Fourier transform of a multiplication of two functions is equivalent to the Fourier transforms of the individual functions convolved together. And vice versa
\begin{equation}
\mathcal{F} \left\{ g \ast h \right\} = \mathcal{F} \left\{ g \right\} \cdot  \mathcal{F} \left\{ h \right\}
\end{equation}
The convolution of the Fourier transform of two functions is equivalent to the individual Fourier transforms of the two functions multiplied together.
Also
\begin{equation}
\mathcal{F} \left\{ \sum\limits_{n=-\infty}^{\infty}  \delta(t-n k) \right\} = \frac{1}{k} \sum\limits_{n=-\infty}^{\infty}  \delta \left( t-\frac{n}{k} \right)
\end{equation}
Using the last three identities, we can write the Fourier transform of \ref{eq:vasconv} as
\begin{equation} 
\mathcal{F} \left\{ p(t) \ast \hat{v}(t) \right\} = \tilde{V}(f) =
\left[ \frac{1}{T_{\Delta}} \cdot P(f) \cdot V(f) \right] \ast
\left[ \sum\limits_{i=-\infty}^{\infty}  \delta \left( t-\frac{i}{T_{\Delta}} \right) \right]
\end{equation}
Where $P(f)$ is the Fourier transform of $p(t)$, $\tilde{V}(f)$ is the Fourier transform of $\tilde{v}(t)$ and $V(f)$ is the Fourier transform of $v(t)$. If we put 
\begin{equation} 
\label{eq:gfdef}
G(f) = \frac{1}{T_{\Delta}} \cdot P(f) \cdot V(f)
\end{equation}
Then we we can arrive to our final conclusion that
\begin{equation}
\label{eq:vspectrum}
\tilde{V}(f) = G(f) \ast
\left[ \sum\limits_{i=-\infty}^{\infty}  \delta \left( t-\frac{i}{T_{\Delta}} \right) \right]
\end{equation}
This simply means that the signal spectrum $G(f)$ repeats at regular intervals throughout the radio spectrum with a frequency of $\frac{1}{T_{\Delta}} = w h f_(f)$.

For example, for a resolution of $800 \cdot 600 @ 60fps$ (ignore blanking intervals for simplicity), we would have a frequency of $\frac{1}{T_{\Delta}} = 800 \cdot 600 \cdot 60 = 28.8 MHz$. This means that we would expect to find such a signal centred at DC, 28.8 MHz, 57.6 MHz, 86.4 MHz, etc.

\subsection{Sampling Rate}

In order to obtain the full video signal, we need to know what is the minimum rate at which we need to sample the baseband. We saw from \ref{eq:vspectrum} that $G(f)$ repeats at the regular intervals. In fact we can see from \ref{eq:gfdef} that it is $G(f)$ that contains the actual video information
$$G(f) = \frac{1}{T_{\Delta}} \cdot P(f) \cdot V(f)$$
As a consequence from the Whittaker--Shannon interpolation formula that we used to construct $v(t)$, we know that $V(f)$ is band limited so that $V(f) = 0$ for all $|f| \geq \frac{1}{2 T_{\Delta}}$. This implies that $G(f)$ is also band limited to the same range (because of the multiplication). Therefore we need to have a receiver with a sampling rate of at least $\frac{1}{T_{\Delta}}$ in order to reconstruct $G(f)$. Note however that $g(t)$ contains the video signal convolved with the shape of the pixel $p(t)$. We can however assume $P(f)=1$ for $|f| \le \frac{1}{2 T_{\Delta}}$ and zero otherwise. Therefore $G(f) = V(f)$ which would imply a sinc shaped $p(t)$. This in fact is not the case since different video cards have different pixel shapes\cite{kuhn2003compromising}. However in practice, this is enough to obtain a good approximation of the signal.

\section{Reception}

\subsection{Theory}
In order to receive a copy of the video signal, we will need to apply a bandpass filter centred at a multiple of the pixel frequency. The width of the bandpass needs to be the same as the band limited video signal we want to eavesdrop. Then we need to generate an internal clock that runs at the pixel rate and use that to obtain the estimated pixel intensities. Luckily these steps are already done in a typical AM (Amplitude Modulation) receiver. Unfortunately the signal bandwidth is much wider than what a typical off-the-shelf AM receiver can handle.

However some Software Defined Radio receivers can provide the required bandwidth. The Ettus Research USRP B200 can provide more than 50 MHz of realtime bandwidth which is more than what is required to eavesdrop video signals. Software Defined Radios provide the radio samples as a quadrature vector which angle relates to the instantaneous phase in relation to the internal hardware oscillator of the hardware. The size of the vector is proportional to the received amplitude for each sample. For AM reception, we do not need the phase information since the video signal is being transmitted with a constant frequency. Therefore we can make a simple AM demodulator by taking the length of the vector of each sample. When our digital sampling rate matches the pixel rate, then each sample will contain an estimate of the pixel intensity at that time.

\subsection{Practice}
The VESA standard specifies an allowed frequency tolerance for $f_{p}$ of 0.5\%. Therefore due to hardware limitations, we might not be able to accurately tune the receiver sampling rate with enough accuracy to match the transmitted pixel rate. This means that some resampling needs to be done in software in order to create a new digital signal that matches as close as possible the transmitted pixel rate in which each sample corresponds to a pixel.

It is also true that we might be able to recover a lot of information from a signal even if our receiver is not capable of sampling the full width of the video spectrum. This means we will only receive a of $g(f)$ with a bandpass filter applied to it. In this case the digital resampling could interpolate the received samples to fit into multiple pixels. Therefore it is useful to introduce a measurement $\kappa$ of how accurate we can reconstruct pixels of $g(t)$. This measurement can be the number of input samples that have gone into constructing a single output pixels. Therefore
\begin{equation}
\kappa = receiver\_samplerate \cdot T_{\Delta} = \frac{receiver\_samplerate}{w h f_{f}}
\end{equation}
If $\kappa \geq 1$, then we can deconstruct individual pixels. If $\kappa < 1$ we are looking at a band-passed version of the signal and consecutive reconstructed pixels will depend on one another. Please refer to Figure \ref{fig:samplerates} for a visual comparison of different values of $\kappa$.

\begin{figure}[p!]
\centering
\begin{subfigure}[b]{0.45\textwidth}
  \includegraphics[width=\textwidth]{sr_original}
  \caption{Transmitted image}
\end{subfigure} ~
\begin{subfigure}[b]{0.45\textwidth}
  \includegraphics[width=\textwidth]{sr_40MHz_at_190MHz}
  \caption{$\kappa = 1.003$}
\end{subfigure} ~
\begin{subfigure}[b]{0.45\textwidth}
  \includegraphics[width=\textwidth]{sr_30MHz_at_190MHz}
  \caption{$\kappa = 0.753$}
\end{subfigure} ~
\begin{subfigure}[b]{0.45\textwidth}
  \includegraphics[width=\textwidth]{sr_20MHz_at_190MHz}
  \caption{$\kappa = 0.502$}
\end{subfigure} ~
\begin{subfigure}[b]{0.45\textwidth}
  \includegraphics[width=\textwidth]{sr_10MHz_at_190MHz}
  \caption{$\kappa = 0.251$}
\end{subfigure} ~
\begin{subfigure}[b]{0.45\textwidth}
  \includegraphics[width=\textwidth]{sr_5MHz_at_190MHz}
  \caption{$\kappa = 0.125$}
\end{subfigure}
\caption{The transmitted image (a) is being sent at a pixel rate of 800x600 at 60 fps with $w = 1056$, $h = 628$ and $f_{f} = 60.11$. Each of the screenshots shows the reconstructed image with the receiver (USRP B2000) running at different sample rates: (b) - 40 MHz, (c) - 30 MHz, (d) - 20 MHz, (e) - 10 MHz and (f) - 5 MHz}
\label{fig:samplerates}
\end{figure}

Another possible source of distortions is the usage of multiple wires to transmit colour information. These signals interfere with each other and the resulting pixel intensity becomes a complicated mix of the signals generated in these wires. Therefore colour information would be very difficult to obtain. Nevertheless the outlined strategy allows an eavesdropper to detect changes in colour since this would result in different signal amplitudes arriving at the receiver for each pixel.

This means that we would not be able to decode colour. Very often we can't even decode grayscale images being displayed either. This is because even grayscale pixels will depend on more than one wire to display the correct shade. What we will see is that each decoded pixel will have an intensity that relates to the colour of the original pixel in a very non-linear and non-unique fashion. Nevertheless practice shows that this is enough to read text from the remote screen as you have seen in the screenshots on Figure \ref{fig:samplerates}.

Dr Markus Kuhn has shown that the outlined procedure even works for decoding digital video signals\cite{kuhn2005electromagnetic}. If the pixel intensities are transmitted as a bit pattern rather than an analogue amplitude, the received intensity will be the average intensity in the bit pattern. The redundancy of bit patterns means that even in such case we can still receive a signal that will allow us to read text on the remote screen. 

\section{Resolution and Framerate Detection}

The outlined method for decoding the video stream relies on us having the exact values for $w$, $h$ and $f_{f}$. This is not very practical for a real world attack. We would require the adversary not to have any knowledge of the target system whatsoever. Therefore we would need to infer all of these parameters remotely by exploiting some typical characteristics of a video signal. Namely the fact that the signal is periodic.

As we have already seen, a video signal consists of video frames with width $w$ and height $h$. The speed at which such frames are transmitted is $f_{f}$. However most of the time any two consecutive frames would be identical. This would be the case if the victim is looking at static text on the screen. Therefore we can regard the transmitted signal as a repeating periodic signal of a single frame. Which means if we analyse the signal for repeating patterns we should be able to spot the repeating video frames in it and use that to estimate $f_{f}$.

\subsection{Introduction to Autocorrelation}
A good way of doing so would be calculating the discrete autocorrelation of the incoming signal which is defined as 
$$ R_{vv}(j)=\sum\limits_{i} v_{i} \bar{v}_{i-j} $$
Where $\bar{v}$ is the complex conjugate of $v$ and $j$ is the \textbf{samples lag}. In our case $v$ is the set of real samples, therefore $\bar{v} = v$. Basically the autocorrelation is a measure of the similarity of a signal with the same signal shifted by a lag. The higher the value of $R_{vv}(j)$ is, the more similar the function is at that lag. Therefore by analysing $R_{vv}(j)$, we would be able to spot any repeating patterns inside $v_{i}$.

Figure \ref{fig:autocorr} shows the discrete autocorrelation of a video signal that was captured using the Mirics USB dongle at 8 MHz (i.e. $8,000,000$ samples per second) sampling rate. The vertical axis is logarithmic, containing the decibels of the autocorrelation i.e. $10 \times log10( R_{vv} )$. The horizontal axis contains the time lag in milliseconds. This was converted from the lag in samples $j$ by $x = j \cdot T_{delta}$.

\begin{figure}[h]
\centering
  \includegraphics[width=0.96\textwidth]{autocorr}
  \caption{Autocorrelation of a signal}
  \label{fig:autocorr}
\end{figure}

This particular example consists of $524,288$ samples which is $65.536$ milliseconds of recording. This allows us to estimate the autocorrelation up to half of it, namely $32.768$ milliseconds. The reason is because $R_{vv}(j)=-R_{vv}(j)$ due to symmetry. We can see a peak at $16.672$ milliseconds which corresponds to a frequency of $f_{f} = \frac{1}{0.016672} = 59.98 Hz$. This is our estimate for $f_{f}$. And it makes sense, $60 Hz$ is a common refresh rate for contemporary video displays.

\subsection{Aliasing}
We should be able to see the peak repeating again at $2 \times 16.672 = 33.344 ms$. This is analogous to comparing every two consecutive frames with every two other frames. However, as this falls just outside our window it would be seen as an alias at $65.536-33.344 = 32.192 ms$. The same applies for the third peak at $3 \times 16.672 = 50.016 ms$ which would be comparing each block of three consecutive frames with the next three. This will alias at $65.536 - 50.016 = 15.52 ms$. This explains the second small peak we see on the left of the main one at $16.672 ms$ but we see that the difference in its power is already several decibels lower. Therefore any further aliasing induced would be even smaller.

It is therefore important to choose the number of samples to be used for the autocorrelation wisely in order to avoid aliasing. We need to have enough full frames in our data so that the autocorrelation can physically work. The minimum frequency we can get from an autocorrelation is $\frac{2 \times receiver\_samplerate}{numer\_of\_samples}$. It could be easily shown that the second alias is
\begin{equation}
a_{2} = \frac{1}{\frac{2}{f_{min}}-\frac{2}{f}} = \frac{1}{\frac{number\_of\_samples}{receiver\_samplerate}-\frac{2}{f}}
\end{equation}
Where $f$ is the frequency we are going to see the alias for. Therefore if we want to be able to pick up a repeating pattern of at least $f_{min} = 50 Hz$, we would need $\frac{2 \times receiver\_samplerate}{50}$ which in our case is $320,000$ samples. However having such a small number of samples if we observe a signal at 85 Hz we would discover a strong alias at $\frac{1}{\frac{320,000}{8,000,000}-\frac{2}{85} } = 60.714$ Hz which would be very misleading and will look like a legitimate signal at $60.714$ Hz. 

Let us assume $f_{lo} \leq f_{f} \leq f_{hi}$. In order to minimise aliasing, we can try to keep the second alias $a_{2}$ outside this range. Therefore we need to pick the number of samples so that
$$ \frac{1}{\frac{number\_of\_samples}{receiver\_samplerate}-\frac{2}{f_{lo}}} \leq f_{lo} \,\,\,\,\ or \,\,\,\,\ \frac{1}{\frac{number\_of\_samples}{receiver\_samplerate}-\frac{2}{f_{hi}}} \geq f_{hi} $$
Which means that
$$ 3 \times \frac{receiver\_samplerate}{f_{lo}} \leq number\_of\_samples $$
or
$$3 \times \frac{receiver\_samplerate}{f_{hi}} \geq number\_of\_samples $$
and don't forget that in the same time we have the condition
$$\frac{2 \times receiver\_samplerate}{f_{lo}} \leq numer\_of\_samples$$
so in the end, in order to eliminate the second alias in the desired region from $f_{lo}$ to $f_{hi}$ in the autocorrelation, we would need to choose the number of samples to be
\begin{equation}
\label{eq:autocorraliasingcondition}
numer\_of\_samples \geq 3 \times \frac{receiver\_samplerate}{f_{lo}}
\end{equation}
To paraphrase this result, we will need to do the autocorrelation on at least 3 frames worth of samples for the lowest frequency in our range $f_{lo}$ to avoid second order aliasing interfering with our plot.

\subsection{Number of Lines in a Frame}

We were able to estimate $f_{f}$ using an autocorrelation and eliminate aliased signals in the frequency window that we are looking to detect repetitions. However, we still have not estimated neither $w$, nor $h$. As a matter of fact we might not be able to measure $w$ at all. The reason is that $w$ might be a continuous signal and could potentially contain any number of pixels. However, each line of $w$ pixels repeat $h$ times in a frame. There are, however, the blanking intervals that repeat on each line. Therefore we can use their repeating nature to estimate $h$ from the autocorrelation plot.

For this reason we need to zoom in to our plot. Refer to Figure \ref{fig:autocorr_zoomed} which is a zoomed out version of the figure we saw in the previous subsection. If we still have our allowed $f_{f}$ window between $f_{lo}$ and $f_{hi}$, then we can assume that $h$ also spans between $h_{lo}$ and $h_{hi}$. We would therefore expect to see a peak between $\frac{1}{f_{hi} \cdot h_{hi}}$ and $\frac{1}{f_{lo} \cdot h_{lo}}$ milliseconds. If we detect a peak at time $t_{peak}$, then this would correspond to $h = round \left( \frac{1}{t_{peak} \cdot f_{f}} \right)$. Note that we should have already estimated $f_{f}$.

\begin{figure}[h]
\centering
  \includegraphics[width=0.96\textwidth]{autocorr_zoomed}
  \caption{Zoomed in version of \ref{fig:autocorr}}
  \label{fig:autocorr_zoomed}
\end{figure}

Our plot shows a peak at $0.017875$ milliseconds. Since we already estimated $f_{f} = 59.98 Hz$, then $h = round \left( \frac{1000}{0.017875 \times 59.98} \right) = 933$ lines. If we look at a list of VESA video modes, we can see that there is a conveniently close video mode that corresponds to $h = 933$ and $f_{f} = 59.98$. It is the ``1440x900 @ 60 Hz'' video mode with $w = 1904$ pixels, $h = 932$ lines and $f_{f}=60 Hz$. Therefore we could use the value of $w = 1904$ for our estimate of the number of pixels in a row. This will keep the aspect ratio correct.

\subsection{Errors}

There are errors that [introduced by sampling and height could be +- 6 for ex]

\chapter{Practical attack} 

Before explaining how the software works internally, let's present a demonstration of a practical video eavesdropping attack. Its main aim is to show the ease with which such an attack could happen. In the meantime this will give an opportunity to explain the characteristics of the received signal and how they are exploited.

As in the real world, we will start with no knowledge of the victim's system. We will estimate the frequency at which the emission strength has the best Signal to Noise Ratio (SNR). We will then analyse the signal to detect the resolution and refresh rate of the screen. We will afterwards lock onto the signal and try to recover the original video. We will also discuss some techniques that could be utilized to improve the quality of the image.

\section{The Setup}

The choice for a SDR front-end for this demonstration is a USRP B200\footnote{Refer to \ref{sec:hw} Hardware for discussion on currently available SDR devices.}. Depending on the particular requirements, an attacker might prefer mobility over accuracy and choose the smaller Mirics FlexiTV\texttrademark MSi3101 SDR USB Dongle. However, for this demonstration we will attempt to obtain the highest possible resolution. Therefore we need an SDR radio that is capable of obtaining a wide bandwidth. This makes the 32 MHz of Bandwidth that the USRP provides a much better choice than the 8 MHz available from the Mirics dongle.

[Antenna]

[Location of victim]

[Victim's fonts]

\section{Preparation}

\chapter{Implementation} 

\section{Hardware}
\label{sec:hw} 

\begin{figure}[h!]
 
  \centering
    \includegraphics[width=0.9\textwidth]{equipment}
    \caption{From left to right: MSi3101, AverMedia antenna and USRP B200}
\end{figure}


\section{Architecture}

\subsection{Main Library}

\subsection{JavaGUI}

\subsubsection{Multiplatform}

\subsubsection{Graphical Interface}

\subsubsection{Visualisation}

\section{Digital Signal Processing}

\subsection{Overview}

\subsection{Signal Reconstruction}

\subsubsection{Multithreading}

\subsubsection{Demodulation}

\subsubsection{Re-sampling}

\subsubsection{Auto Gain}

\subsubsection{Low-pass}

\subsection{Frame Synchronization}

\subsection{Resolution Detection}

\subsection{Extended Bandwidth}

\subsection{Testing and Benchmarking}

\chapter{Summary and Conclusions} 

[Summarize results]

[potential future work]

\appendix
\singlespacing

\bibliographystyle{unsrt} 
\bibliography{dissertation} 

\end{document}
